% 若编译失败，且生成 .synctex(busy) 辅助文件，可能有两个原因：
% 1. 需要插入的图片不存在：Ctrl + F 搜索 'figure' 将这些代码注释/删除掉即可
% 2. 路径/文件名含中文或空格：更改路径/文件名即可

% --------------------- 文章宏包及相关设置 --------------------- %
% >> ------------------ 文章宏包及相关设置 ------------------ << %
% 设定文章类型与编码格式
\documentclass[UTF8]{article}		
\input{../.config/config_for_MCPExperiment.tex}



%%%%%%%%%%%%%%%%%%%%%%%%%%%%%%%%%%%%%%%%%%%%%%%%%%%%%%%%%%%%%%%%
% 仅需修改页眉、实验名称、实验日期
%%%%%%%%%%%%%%%%%%%%%%%%%%%%%%%%%%%%%%%%%%%%%%%%%%%%%%%%%%%%%%%%


%%%%%%%%%%%%%%%%%% 1. 修改页眉内容 %%%%%%%%%%%%%%%%%%
\rhead{MCPE-08 表中正负零统计与数据排序 (2025.11.26, 徐睿琳)}

\begin{document}
\begin{center}\large
    \vspace*{-0.8cm}
    \noindent{\huge\bfseries《\ \ 微\ \ 机\ \ 原\ \ 理\ \ 实\ \ 验\ \ \ 》\ \ 实\ \ 验\ \ 报\ \ 告 }
    \\\vspace{0.1cm}
    \noindent{
    {\bfseries 
%
%%%%%%%%%%%%%%%%%% 2. 修改实验名称 %%%%%%%%%%%%%%%%%%
    实验名称：\uline{\hspace{2.2cm} 表中正负零统计与数据排序 \hspace{2.2cm}}
%
    }\hspace{0.4cm}
    指导教师：\uline{\hspace{0.8cm}肖山林     \hspace{0.8cm}}
    }
    \\\vspace{0.1cm}
    \noindent
    {
    姓名：\uline{\,\,\,徐睿琳\,\,\,}\hspace{0.2cm}
    学号：\uline{\,\,\,{ 23342107}\,\,\,}\hspace{0.2cm}
    专业/班级：\uline{\,\,\,{微电子/三班}\,\,\,}\hspace{0.2cm}
    分组序号：\uline{\,\,\,{A412-A05}\,\,\,}
    }
    \\\vspace{0.1cm}
    \noindent{
%
%%%%%%%%%%%%%%%%%% 3. 修改实验日期 %%%%%%%%%%%%%%%%%%
    实验日期：\uline{\,\,{2025.11.26}\,\,}\hspace{0.2cm}
%
    实验地点：\uline{\,\,\,教学楼{ A412}\,\,\,}\hspace{0.2cm}
    是否调课/补课：\uline{\hspace{0.26cm}否 \hspace{0.26cm}}\hspace{0.2cm}
    成绩：\uline{\hspace{2cm}}
    }
\end{center}
\vspace{-0.4cm}
\noindent\rule{\textwidth}{0.075em}   % 分割线
\vspace{-1.0cm}

% 生成目录
\setcounter{tocdepth}{3}  % 目录深度为 2（不显示 subsubsection）
\noindent\tableofcontents\thispagestyle{fancy}   % 显示页码、页眉等

% ------------------------ 文章信息区 ------------------------ %
% ------------------------ 文章信息区 ------------------------ %



%%%%%%%%%%%%%%%%%%%%%%%%%%%%%%%%%%%%%%%%%%%%%%%%%%%%%%%%%%%%%%%%%%%%%%%%%%%%%%%%%
%%%%%%%%%%%%%%%%%%%%%%%%%%%%%%%%% 下面是正文内容 %%%%%%%%%%%%%%%%%%%%%%%%%%%%%%%%%
%%%%%%%%%%%%%%%%%%%%%%%%%%%%%%%%%%%%%%%%%%%%%%%%%%%%%%%%%%%%%%%%%%%%%%%%%%%%%%%%%

\section{实验内容与设计}

\subsection{数据排列实验}

\subsubsection{流程图}

\begin{figure}[H]\centering
	\includegraphics[width=1\columnwidth]{MCPE-07 Data Arrangement Experiment/source/pic/外循环流程图.png}
	\caption{外循环流程图}
\end{figure}

\begin{figure}[H]\centering
	\includegraphics[width=1\columnwidth]{MCPE-07 Data Arrangement Experiment/source/pic/内循环流程图.png}
	\caption{内循环流程图}
\end{figure}

\subsubsection{代码}

\begin{lstlisting}[language={[x86masm]Assembler}]
CODE    SEGMENT
ASSUME  CS:CODE

START:
MOV     AX, CS
MOV     DS, AX


MOV     BYTE PTR [0200H], 25
MOV     BYTE PTR [0201H], 46
MOV     BYTE PTR [0202H], 3
MOV     BYTE PTR [0203H], 75
MOV     BYTE PTR [0204H], 5
MOV     BYTE PTR [0205H], 30

; 调用子程序
CALL    BUBBLE_SORT

JMP     $

BUBBLE_SORT PROC NEAR
MOV     CX, 6                   
CMP     CX, 1                   ; 如果长度<=1
JBE     SORT_DONE

DEC     CX                      ; 外层循环次数 = 长度-1

OUTER_LOOP:
MOV     BX, 0                   
MOV     SI, 0                   

INNER_LOOP:
; 检查 SI+1 是否越界
MOV     AX, SI
INC     AX
CMP     AX, 6                   ; 比较索引是否超出数组边界
JAE     END_INNER_LOOP          ; 如果SI+1 >= 数组长度，结束内层循环

; 相邻元素比较
MOV     AL, [0200H + SI]        ; AL = 当前
MOV     AH, [0200H + SI + 1]    ; AH = 下一个

CMP     AL, AH
JBE     NO_SWAP                 ; AL <= AH，不交换

; 交换
XCHG    AL, AH
MOV     [0200H + SI], AL        ; 存回当前元素
MOV     [0200H + SI + 1], AH    ; 存回下一个元素
MOV     BX, 1                   ; 设置交换标志

NO_SWAP:
INC     SI                      ; 移动到下一对元素
JMP     INNER_LOOP

END_INNER_LOOP:
CMP     BX, 0
JE      SORT_DONE               ; 整个内循环没发生交换，排序完成

LOOP    OUTER_LOOP              ; 继续下一轮排序

SORT_DONE:
RET
BUBBLE_SORT ENDP

CODE    ENDS
END     START
\end{lstlisting}

\subsection{统计表中正负数，0}

\subsubsection{流程图}

\begin{figure}[H]\centering
	\includegraphics[width=1\columnwidth]{MCPE-08 Count the number of positive, negative, and zero values in the table/source/pic/流程图.png}
	\caption{流程图}
\end{figure}

\subsubsection{代码}

\begin{lstlisting}[language={[x86masm]Assembler}]
DATA    SEGMENT
TABLE   DW  20                      ; 表长
DW  -10,-9,-8,-7,-6,-5,-4,-3,-2,-1
DW  0,1,2,3,4,5,6,7,8,9     

NEGNUM  DW  0       ; -
POSNUM  DW  0       ; + 
ZERONUM DW  0       ; 0
DATA    ENDS

CODE    SEGMENT
ASSUME  CS:CODE, DS:DATA

START:

MOV     AX, DATA
MOV     DS, AX

; 初始化计数器
XOR     BX, BX          ; -
XOR     DX, DX          ; 0
XOR     AX, AX          ; +

MOV     CX, TABLE       
MOV     SI, 2           ; 指向第一个数据（跳过一个字）

COUNT_LOOP:
; 读取当前数字并分类
CMP     WORD PTR [TABLE+SI], 0
JL      IS_NEG          ; Jump if Less SF ≠ OF
JG      IS_POS          ; Jump if Greater ZF=0 AND SF=OF

; 前两个都没跳转就是0
INC     DX              
JMP     NEXT

IS_NEG:
INC     BX              
JMP     NEXT

IS_POS:
INC     AX              

NEXT:
ADD     SI, 2           ; 移动到下一个字
LOOP    COUNT_LOOP      

MOV     NEGNUM, BX      
MOV     ZERONUM, DX     
MOV     POSNUM, AX      

JMP     $

CODE    ENDS
END     START
\end{lstlisting}

\section{实验结果}

\subsection{数据排列实验}

\begin{figure}[H]\centering
	\includegraphics[width=1\columnwidth]{MCPE-08 Count the number of positive, negative, and zero values in the table/source/pic/字符串初始化.png}
	\caption{排序数值串初始化}
\end{figure}

\begin{figure}[H]\centering
	\includegraphics[width=1\columnwidth]{MCPE-08 Count the number of positive, negative, and zero values in the table/source/pic/第一次内循环.png}
	\caption{第一次内循环}
\end{figure}

\begin{figure}[H]\centering
	\includegraphics[width=1\columnwidth]{MCPE-08 Count the number of positive, negative, and zero values in the table/source/pic/最终排序结果.png}
	\caption{最终排序结果}
\end{figure}

\subsection{统计表中正负数，0}

\begin{figure}[H]\centering
	\includegraphics[width=1\columnwidth]{MCPE-08 Count the number of positive, negative, and zero values in the table/source/pic/表初始化.png}
	\caption{表数据初始化}
\end{figure}

\vspace*{-3mm}
\begin{bluebox}
    初始化表格，第一位是表内数据总数，即20个数据，后面依次是-10...-1,0,1...9，一共10个负数，1个0,9个正数。
\end{bluebox}

\begin{figure}[H]\centering
	\includegraphics[width=1\columnwidth]{MCPE-08 Count the number of positive, negative, and zero values in the table/source/pic/MOVCX, TABLE.png}
	\caption{\texttt{MOV     CX, TABLE}}
\end{figure}

\vspace*{-3mm}
\begin{bluebox}
	将总数设置为循环次数（检查次数）
\end{bluebox}

\begin{figure}[H]\centering
	\includegraphics[width=1\columnwidth]{MCPE-08 Count the number of positive, negative, and zero values in the table/source/pic/-20判断.png}
	\caption{判断 -10 的正负}
\end{figure}

\vspace*{-3mm}
\begin{bluebox}
	使用\texttt{-10-0}，结果为负，触发\texttt{JL}的跳转条件，即\texttt{SF}不等于\texttt{OF}。
\end{bluebox}

\begin{figure}[H]\centering
	\includegraphics[width=1\columnwidth]{MCPE-08 Count the number of positive, negative, and zero values in the table/source/pic/1判断.png}
	\caption{判断 1 的正负}
\end{figure}

\vspace*{-3mm}
\begin{bluebox}
	使用\texttt{1-0}，结果为正，触发\texttt{JG}的跳转条件，即\texttt{SF}等于\texttt{OF}，且ZF不为0。
\end{bluebox}

\begin{figure}[H]\centering
	\includegraphics[width=1\columnwidth]{MCPE-08 Count the number of positive, negative, and zero values in the table/source/pic/结果存放.png}
	\caption{结果存放}
\end{figure}

\vspace*{-3mm}
\begin{bluebox}
	将最终结果存放在指定存储位置。
\end{bluebox}

\section{思考与总结}





























\subsection{实验总结}

本次实验完成了两个汇编程序的设计与实现：

\begin{enumerate}
    \item \textbf{冒泡排序算法}：采用双层循环结构，外层控制排序轮数，内层进行相邻元素比较与交换。通过设置交换标志\texttt{BX}优化算法，当某轮无交换发生时提前结束排序。
    
    \item \textbf{正负零统计}：利用\texttt{CMP}指令与条件跳转指令（\texttt{JL}、\texttt{JG}）实现数据分类。\texttt{JL}判断负数（\texttt{SF≠OF}），\texttt{JG}判断正数（\texttt{ZF=0}且\texttt{SF=OF}），两者均不满足则为零。
\end{enumerate}

\subsection{思考}

\begin{enumerate}
    \item \textbf{有符号数与无符号数的区别}：本实验使用\texttt{JL/JG}处理有符号数，若处理无符号数应使用\texttt{JB/JA}指令。
    
    \item \textbf{寻址方式}：实验中使用\texttt{[TABLE+SI]}的变址寻址方式遍历数组，需注意字类型数据每次偏移量为2。
    
    \item \textbf{标志位的理解}：条件跳转依赖标志寄存器状态，深入理解\texttt{SF}、\texttt{OF}、\texttt{ZF}等标志位对编写正确的分支程序至关重要。
    
    \item \textbf{程序优化}：冒泡排序中引入交换标志可减少不必要的循环次数，体现了算法优化思想。
\end{enumerate}



















\newpage
% 附录
\section*{附录 A\hspace*{20pt} 原始数据记录表}
\addcontentsline{toc}{section}{附录 A\hspace*{6pt} 原始数据记录表} 
\thispagestyle{fancy} 
\begin{figure}[H]\centering
    %\includegraphics[width=\columnwidth]{LCE-04-场效应管/assets/appendix/预习报告.png}
\end{figure}

% 附录
\newpage
\vspace*{\fill}\begin{center}\Huge{\bfseries 
    附录 B\hspace*{20pt} 实验预习报告
}\end{center}\addcontentsline{toc}{section}{附录 B\hspace*{6pt} 实验预习报告} 
\begin{figure}[H]\centering
    %\includegraphics[width=\columnwidth]{LCE-04-场效应管/assets/appendix/预习报告.png}
\end{figure}
\vspace*{\fill}
\thispagestyle{fancy} 
%\includepdf[pages={-}]{LCE-04-场效应管/preview/LCE-04 (preview report).pdf}


% 附录
\section*{附录 C \hspace*{20pt} asm Codes}
\addcontentsline{toc}{section}{附录 C \hspace*{6pt} asm Codes} 
\thispagestyle{fancy} 
%\lstinputlisting{d:/a_RemoteRepo/GH.MatlabCodes/本科课程代码/基础物理实验/Ex_9/Ex_9_analysis_mfile.m}



\end{document}

% VScode 常用快捷键：

% F2:                       变量重命名
% Ctrl + Enter:             行中换行
% Alt + up/down:            上下移行
% 鼠标中键 + 移动:           快速多光标
% Shift + Alt + up/down:    上下复制
% Ctrl + left/right:        左右跳单词
% Ctrl + Backspace/Delete:  左右删单词    
% Shift + Delete:           删除此行
% Ctrl + J:                 打开 VScode 下栏(输出栏)
% Ctrl + B:                 打开 VScode 左栏(目录栏)
% Ctrl + `:                 打开 VScode 终端栏
% Ctrl + 0:                 定位文件
% Ctrl + Tab:               切换已打开的文件(切标签)
% Ctrl + Shift + P:         打开全局命令(设置)

% Latex 常用快捷键：

% Ctrl + Alt + J:           由代码定位到PDF


